% ==STYLE-BEGIN==
\documentclass[UTF8,11pt]{ctexart}
\usepackage{amsmath,amssymb,amsthm}
\usepackage[a4paper,margin=2.4cm,top=2.3cm,bottom=2.3cm]{geometry}
\usepackage{fancyhdr}
\usepackage{enumitem}
\usepackage{hyperref}
\hypersetup{hidelinks}

% ---- 版式 ----
\linespread{1.15}
\setlength{\parskip}{0.35em}
\setlength{\parindent}{1.6em}
\setlist{nosep,leftmargin=2em}

% ---- 页眉 ----
\pagestyle{fancy}
\fancyhf{}
\fancyhead[L]{\small\coursename}
\fancyhead[R]{\small\thepage}
\renewcommand{\headrulewidth}{0.4pt}

% ---- 定理环境：中文用正体，标题加粗，按小节共用一条流水号 ----
\newtheoremstyle{cnthm}{0.7em}{0.7em}{\normalfont}{}{\bfseries}{}{0.5em}{}
\theoremstyle{cnthm}
\newtheorem{thm}{定理}[section]
\newtheorem{lem}[thm]{引理}
\newtheorem{prop}[thm]{命题}
\newtheorem{cor}[thm]{推论}
\newtheorem{defn}[thm]{定义}
\newtheorem{rem}[thm]{注}
\newtheorem{exa}[thm]{例}
\newtheorem{conj}[thm]{猜想}
\renewcommand{\proofname}{\bfseries 证明}
\providecommand{\coursename}{}
% ==STYLE-END==
\renewcommand{\coursename}{高等代数\,·\,1116}
\xeCJKDeclareCharClass{CJK}{"2150 -> "218F}
\title{高等代数 1116}
\date{}
\begin{document}
\maketitle

\section{回顾: 方程组与三套语言}

\subsection{线性方程组与三个问题}

\begin{equation*}
\begin{cases}
a_{11}x_1+a_{12}x_2+\cdots+a_{1n}x_n=b_1,\\
a_{21}x_1+a_{22}x_2+\cdots+a_{2n}x_n=b_2,\\
\qquad\qquad\cdots\cdots\\
a_{m1}x_1+a_{m2}x_2+\cdots+a_{mn}x_n=b_m,
\end{cases}
\tag{$*$}
\end{equation*}

\textbf{Q}: (1) 方程组是否有解? (2) 如何求解? (3) 解集的结构如何?

\textbf{A}: (1) 阶梯形方程组中没有矛盾方程; (2) Gauss 消元法——用自由未知量表示主未知量.

\subsection{语言一: $m$ 维向量空间 $K^m$}

$(*)$ $\iff$ 向量方程
\[
x_1\alpha_1+x_2\alpha_2+\cdots+x_n\alpha_n=\beta .
\]
由此可以把解析几何里的概念全部引进来: 线性组合、线性表示、线性相关/无关; 极大无关组、秩; 子空间及其基与维数.
\[
(*)\ \text{有解}\iff\beta\ \text{可被}\ \alpha_1,\cdots,\alpha_n\ \text{线性表示}
\iff \beta\in L(\alpha_1,\cdots,\alpha_n).
\]

若 $\beta=\mathbf 0$: 取 $\alpha_1,\cdots,\alpha_n$ 的极大无关组, 再把多余的向量用极大无关组表示; \textbf{每一个表示就给出一个解向量}, 这些解向量构成一个\textbf{基础解系}.

\subsection{语言二: 矩阵}

增广矩阵 $\bar A$ 经初等行变换化为(简化)行阶梯形矩阵, 于是
\[
(*)\ \text{有解}\iff\text{最后一列无主元}\iff r(A)=r(\bar A).
\]
若 $\beta=\mathbf 0$, 则
\[
(*)\ \text{有非零解}\iff r(A)<n .
\]

\section{矩阵代数}

下面把"矩阵"本身当作运算对象来研究——这是继方程组、向量之后的第三套语言(三套语言中的第三套).

\begin{defn}[矩阵]
数域 $K$ 上 $mn$ 个数排成的一个 $m$ 行 $n$ 列的数表
\[
\begin{pmatrix}
a_{11}&a_{12}&\cdots&a_{1n}\\
a_{21}&a_{22}&\cdots&a_{2n}\\
\vdots&\vdots&&\vdots\\
a_{m1}&a_{m2}&\cdots&a_{mn}
\end{pmatrix}
\]
称为 $K$ 上的一个 $m\times n$ 矩阵. $K$ 上全体 $m\times n$ 矩阵记为 $M_{m,n}(K)$.
\end{defn}

\begin{rem}
(1) 矩阵一般用 $A,B,C,P,Q$ 表示; 元素用 $a,b,c$ 表示. 例如 $A$ 的第 $i$ 行与第 $j$ 列交叉处的元素用 $a_{ij}$ 表示, 也可用 $A(i,j)$;

(2) 若 $A$ 的元素全是 $0$, 则称 $A$ 是 $O$ 矩阵, 记作 $O$;

(3) 如果 $A\in M_{m,n}(K)$ 且 $m=n$, 则称 $A$ 为 $n$ 阶方阵, 全体 $n$ 阶方阵记为 $M_n(K)$.
\end{rem}

\section{矩阵的加法与数乘}

设 $A=(a_{ij})\in M_{m,n}(K)$, $B=(b_{ij})\in M_{m,n}(K)$, $k\in K$.

\begin{defn}[加法与数乘]
(1) \textbf{加法}: $A+B\triangleq(a_{ij}+b_{ij})\in M_{m,n}(K)$;

(2) \textbf{数乘}: $kA\triangleq(ka_{ij})\in M_{m,n}(K)$.
\end{defn}

\begin{thm}
$(M_{m,n}(K),+,\cdot)$ 是一个线性空间, 其中
\[
-A\triangleq(-a_{ij}),\qquad A-B\triangleq A+(-B).
\]
\end{thm}

\textbf{同构}　把 $A=(a_{ij})$ 的 $mn$ 个元素按顺序排成一列, 就得到 $K^{mn}$ 中的一个向量, 且加法、数乘一一对应, 故
\[
(M_{m,n}(K),+,\cdot)\cong(K^{mn},+,\cdot).
\]

\section{矩阵的乘法}

\textbf{动机}　线性方程组可以写成"$ax=b$"的形状:
\[
\begin{cases}
a_{11}x_1+\cdots+a_{1n}x_n=b_1,\\
\qquad\cdots\cdots\\
a_{m1}x_1+\cdots+a_{mn}x_n=b_m
\end{cases}
\iff
AX=\beta,
\]
其中
\[
A=(a_{ij}),\qquad X=\begin{pmatrix}x_1\\\vdots\\x_n\end{pmatrix},\qquad
\beta=\begin{pmatrix}b_1\\\vdots\\b_m\end{pmatrix}.
\]
为了让这个写法成立, 需要定义矩阵与向量的乘法.

\subsection{矩阵与向量相乘}

\begin{defn}
设 $A=(a_{ij})\in M_{m,n}(K)$, $X=(x_1,x_2,\cdots,x_n)^{\mathrm T}\in M_{n,1}(K)$, 定义
\[
AX=
\begin{pmatrix}
a_{11}&a_{12}&\cdots&a_{1n}\\
a_{21}&a_{22}&\cdots&a_{2n}\\
\vdots&\vdots&&\vdots\\
a_{m1}&a_{m2}&\cdots&a_{mn}
\end{pmatrix}
\begin{pmatrix}x_1\\x_2\\\vdots\\x_n\end{pmatrix}
\triangleq
\begin{pmatrix}
a_{11}x_1+a_{12}x_2+\cdots+a_{1n}x_n\\
a_{21}x_1+a_{22}x_2+\cdots+a_{2n}x_n\\
\vdots\\
a_{m1}x_1+a_{m2}x_2+\cdots+a_{mn}x_n
\end{pmatrix}.
\]
\end{defn}

\textbf{条件}: $A$ 的列数 $=$ $X$ 的行数.

\subsection{矩阵与矩阵相乘}

\begin{defn}
设 $A\in M_{m,n}(K)$, $B\in M_{n,s}(K)$. 把 $B$ 按列分块 $B=(B_1,B_2,\cdots,B_s)$, 定义
\[
AB=A(B_1,B_2,\cdots,B_s)\triangleq(AB_1,AB_2,\cdots,AB_s).
\]
\end{defn}

\textbf{条件}: $A$ 的列数 $=$ $B$ 的行数.

若 $A=(a_{ij})\in M_{m,n}(K)$, $B=(b_{ij})\in M_{n,s}(K)$, 则 $AB\in M_{m,s}(K)$, 且
\[
AB=
\begin{pmatrix}
\sum\limits_{k=1}^{n}a_{1k}b_{k1}&\cdots&\sum\limits_{k=1}^{n}a_{1k}b_{ks}\\
\vdots&&\vdots\\
\sum\limits_{k=1}^{n}a_{mk}b_{k1}&\cdots&\sum\limits_{k=1}^{n}a_{mk}b_{ks}
\end{pmatrix}.
\]
单独看每个元素: \textbf{$AB$ 的 $(i,j)$ 元 $=$ $A$ 的第 $i$ 行与 $B$ 的第 $j$ 列对应相乘再相加}:
\[
(AB)(i,j)=(a_{i1},a_{i2},\cdots,a_{in})
\begin{pmatrix}b_{1j}\\b_{2j}\\\vdots\\b_{nj}\end{pmatrix}
=\sum_{k=1}^{n}a_{ik}b_{kj}.
\]

\section{矩阵乘法的性质}

\begin{thm}
(1) $(AB)C=A(BC)$(结合律);

(2) $A(B+C)=AB+AC$(左分配律);

(3) $k(AB)=(kA)B=A(kB)$;

(4) $(AB)^{\mathrm T}=B^{\mathrm T}A^{\mathrm T}$, $(A+B)^{\mathrm T}=A^{\mathrm T}+B^{\mathrm T}$.
\end{thm}

\begin{proof}[(1) 结合律]
设 $A\in M_{m,n}(K)$, $B\in M_{n,s}(K)$, $C\in M_{s,t}(K)$.
\[
\bigl((AB)C\bigr)(i,j)=\sum_{k=1}^{s}(AB)(i,k)\,C(k,j)
=\sum_{k=1}^{s}\Bigl(\sum_{l=1}^{n}A(i,l)B(l,k)\Bigr)C(k,j)
=\sum_{k=1}^{s}\sum_{l=1}^{n}A(i,l)B(l,k)C(k,j),
\]
\[
\bigl(A(BC)\bigr)(i,j)=\sum_{l=1}^{n}A(i,l)(BC)(l,j)
=\sum_{l=1}^{n}A(i,l)\Bigl(\sum_{k=1}^{s}B(l,k)C(k,j)\Bigr)
=\sum_{l=1}^{n}\sum_{k=1}^{s}A(i,l)B(l,k)C(k,j).
\]
两者只是\textbf{求和次序不同}(有限和可交换), 故相等.
\end{proof}

\begin{proof}[(4) 转置]
由定义
\[
(AB)^{\mathrm T}(i,j)=(AB)(j,i)=\sum_{k=1}^{n}A(j,k)B(k,i),
\]
\[
(B^{\mathrm T}A^{\mathrm T})(i,j)=\sum_{k=1}^{n}B^{\mathrm T}(i,k)A^{\mathrm T}(k,j)=\sum_{k=1}^{n}B(k,i)A(j,k).
\]
两者相等(数的乘法可交换), 故 $(AB)^{\mathrm T}=B^{\mathrm T}A^{\mathrm T}$.
\end{proof}

\subsection{注一: 矩阵乘法没有交换律}

$AB\neq BA$ 是常态, 分三种情形:

\begin{enumerate}
\item $AB$ 与 $BA$ \textbf{不一定同时有意义}(要求 $A$ 的列数 $=$ $B$ 的行数, 同时 $B$ 的列数 $=$ $A$ 的行数, 一般不成立);
\item 即使都有意义($A\in M_{m,n}(K)$, $B\in M_{n,m}(K)$), 当 $m\neq n$ 时 $AB$ 是 $m$ 阶、$BA$ 是 $n$ 阶, \textbf{阶数都不同};
\item 即使 $m=n$, 也一般不相等. 例如
\[
A=\begin{pmatrix}1&-1\\0&2\end{pmatrix},\qquad
B=\begin{pmatrix}-1&1\\2&3\end{pmatrix},
\]
\[
AB=\begin{pmatrix}-3&-2\\4&6\end{pmatrix},\qquad
BA=\begin{pmatrix}-1&3\\2&4\end{pmatrix},\qquad AB\neq BA .
\]
\end{enumerate}

\subsection{注二: 矩阵有零因子}

\[
AB=O\ \nRightarrow\ A=O\ \text{或}\ B=O .
\]
例如
\[
A=\begin{pmatrix}1&1\\1&1\end{pmatrix},\qquad
B=\begin{pmatrix}1&-1\\-1&1\end{pmatrix},\qquad
AB=\begin{pmatrix}0&0\\0&0\end{pmatrix}.
\]
对照:

\begin{center}
\begin{tabular}{lll}
\hline
对象 & 有没有零因子 & 乘法可交换吗\\
\hline
数域 $K$ & 没有($ab=0\Rightarrow a=0$ 或 $b=0$) & 可交换\\
多项式 & 没有($fg=0\Rightarrow f=0$ 或 $g=0$) & 可交换\\
\textbf{矩阵} & \textbf{有} & \textbf{不可交换}\\
函数(复合) & 有 & 不可交换($f\circ g\neq g\circ f$)\\
\hline
\end{tabular}
\end{center}

所以矩阵乘法只有"结合律 $+$ 分配律 $+$ 与数乘相容", 没有交换律、也没有消去律——凡是用到交换律、消去律的推理, 在矩阵里都要重新检查.

\end{document}
